Code structure by Virginie Van den Schriek

%----------------------------------------------------------------------------------------
%	PACKAGES AND OTHER DOCUMENT CONFIGURATIONS
%----------------------------------------------------------------------------------------
\documentclass[a4paper,10pt,final,fleqn]{article}
\usepackage[frenchb]{babel}
\usepackage{fontenc}
\usepackage{fancyhdr} % Required for custom headers
\usepackage{lastpage} % Required to determine the last page for the footer
\usepackage{extramarks} % Required for headers and footers
\usepackage{graphicx} % Required to insert images
\usepackage[utf8]{inputenc}
\usepackage{apacite}
\usepackage{url}
\usepackage[normalem]{ulem}
\usepackage{enumitem}
\usepackage{mathtools}


% Margins
%\topmargin=-0.45in
\evensidemargin=0in
\oddsidemargin=0in
\textwidth=6in
\textheight=9.0in
\headsep=0.25in 

%\linespread{1.1} % Line spacing

% Set up the header and footer
\pagestyle{fancy}
\lhead{\AuthorName} % Top left header
\chead{\Course} % Top center header
\rhead{\Date} % Top right header
\lfoot{\lastxmark} % Bottom left footer
\cfoot{} % Bottom center footer
\rfoot{\ \thepage\ / 	\pageref{LastPage}} % Bottom right footer
\renewcommand\headrulewidth{0.4pt} % Size of the header rule
\renewcommand\footrulewidth{0.4pt} % Size of the footer rule



%\setlength\parindent{0pt} % Removes all indentation from 
\setcounter{tocdepth}{2}


\newcommand{\Course}{EPHEC} % Course/class
\newcommand{\AuthorName}{Groupe 2.x} % Your name
\newcommand{\Date}

\title{
\parbox{15cm}
{\includegraphics[width=4cm]{ephec.png} \\ 
  Louvain-la-neuve\\
    \vspace{3cm}
	\begin{center}\sf\bfseries\Huge
		\rule{15cm}{1pt}
		\medskip
		Réponses aux questionnaire-guide \\
		\huge APP2 : VoIP(1)\\
		\vspace{.5cm}
		%\Large Enoncé destiné aux 2TL1\\ [-2mm]
		\rule{15cm}{1pt}
	\end{center}
	\vspace{3cm}
}} 

\author{Ramiro-Gonzalez, De Mets, Vergotte, Goerges, Monroe} \date{\today}


\begin{document}
\thispagestyle{empty}

\maketitle

\newpage

\section{Recherche Théorique}
	
	\subsection{Transmissions vocales}
		
		\begin{description}[style=nextline]

			\item[Comment la voix est-elle numérisée? Indiquer les différentes
			étapes. Quels paramètres vont influencer la qualité de la communication?
			Quel débit obtient-on, en fonction de différents cas de figure (téléphonie,
			MP3 vs Hi-Fi, mono vs stéréo, etc )? Expliquer les calculs.]
			Numérisation d'un signal audio : 
				begin{itemize}
					\item La voix constitue un signal \textbf{analogique}.
					\item Ce signal est ensuite échantillonné selon Shannon-Nyquist. Par exemple
					pour un signal de qualité téléphonique, la bande-passante du signal est de 4KHz,
					donc la fréquence d'échantillonnage doit être de 2*4000Hz(Shannon-Nyquist).
					On a donc un taux de 8000 échantillons/seconde.
					\item Ensuite le signal échantillonné est quantifié pour attribuer un code
					binaire à chaque échantillon, typiquement on quantifiera ce signal sur 8 bits,
					ce qui offre \begin{equation} 2^8 = 256 \end{equation} valeurs possibles.
					\item On obtient donc un signal sonore numérisé d'un débit binaire de :
						Taux d'échantillonnage * Codage = \begin{equation} 8000 \ast 8 = 64Kbps \end{equation}.
				\end{itemize}

		\end{description}
\end{document}
